\section{Galois representations and adjoint vanishing}
\label{av-section}

A theorem of Thorne gives a substantial simplification of the
complementary-parity argument. It proves the adjoint vanishing
assertion at each of the automorphic points occurring in the
submitted construction. We give the comparison of representations,
polarizations, coefficient fields, and local conditions explicitly.
The theorem is a classical input; its Taylor--Wiles proof is not
reproduced here. The argument below does not assert that every
deformation point is automorphic, and gives no bound for the length
of an integral congruence module.

\subsection{Local conditions and finite coefficient fields}

Let $p$ be a prime, let $F$ be a number field, and let $E$ be a
finite extension of $\mathbb Q_p$. Write $G_F$ for the absolute
Galois group of $F$. If $W$ is a finite-dimensional continuous
$E$-representation of $G_F$, a continuous one-cocycle is a map
$z:G_F\longrightarrow W$ satisfying
\[
                    z(gh)=z(g)+g z(h).                         \tag{AV1}
\]
A coboundary has the form $z(g)=gX-X$ for one $X\in W$.
The quotient of cocycles by coboundaries is $H^1(F,W)$.
For each finite place $v$, fix a decomposition group $G_v$
and its inertia group $I_v$. Put
\[
 H^1_f(F_v,W)=
 \begin{cases}
 \ker\bigl(H^1(G_v,W)\longrightarrow H^1(I_v,W)\bigr),&v\nmid p,\\
 \ker\bigl(H^1(G_v,W)\longrightarrow
       H^1(G_v,W\otimes_{\mathbb Q_p}B_{\mathrm{cris}})\bigr),&v\mid p.
 \end{cases}                                                   \tag{AV2}
\]
Here $B_{\mathrm{cris}}$ is Fontaine's crystalline period ring,
with its usual continuous local Galois action; the tensor product
has the diagonal action. Define
\[
 H^1_f(F,W)=\{\alpha\in H^1(F,W):
                    \alpha|_{G_v}\in H^1_f(F_v,W)\text{ for all finite }v\}.
                                                               \tag{AV3}
\]
Changing a decomposition group by conjugation does not change
this condition. Indeed (AV1) gives
\[
 z(g)=\gamma z(\gamma^{-1}g\gamma)+z(\gamma)-g z(\gamma).
                                                               \tag{AV4}
\]
Thus the cocycles on conjugate local groups differ, after the
natural conjugation isomorphism, by a coboundary. The same
calculation applies after tensoring with $B_{\mathrm{cris}}$.

We shall also use representations and cocycles with values in
$\overline{\mathbb Q}_p$, with its $p$-adic topology. The following
elementary point avoids an implicit assumption on their field
of definition.

\begin{lemma}\label{av-coefficients}
Every compact subgroup of
$\mathrm{GL}_d(\overline{\mathbb Q}_p)$ is contained in
$\mathrm{GL}_d(E)$ for some finite extension
$E/\mathbb Q_p$. Consequently a continuous representation of a
profinite group and one continuous cocycle for it are defined,
in their given bases, over a common finite extension of
$\mathbb Q_p$.
\end{lemma}
\begin{proof}
There are countably many finite subfields of
$\overline{\mathbb Q}_p$ over $\mathbb Q_p$. Here is a useful
justification. If $K=\mathbb Q_p(\alpha)$ is finite, approximate
the coefficients of the monic minimal polynomial of $\alpha$
by rational numbers, keeping the degree fixed. Hensel's lemma
in the complete field $K$ gives a root $\beta\in K$ arbitrarily
close to $\alpha$ of one of the resulting polynomials in
$\mathbb Q[X]$. Take
\[
 |\alpha-\beta|_p<
        \min_{\alpha'\ne\alpha}|\alpha-\alpha'|_p,
                                                               \tag{AV5}
\]
where $\alpha'$ runs over the other $\mathbb Q_p$-conjugates.
Every automorphism fixing $\beta$ preserves the $p$-adic
absolute value, so (AV5) forces it to fix $\alpha$.
Thus $\mathbb Q_p(\alpha)\subseteq\mathbb Q_p(\beta)$,
and the reverse inclusion follows from $\beta\in K$.
The degree-one case is immediate. There are only countably
many such algebraic numbers $\beta$, which proves the assertion.

Enumerate these finite subfields as $E_1,E_2,\ldots$, and let
$H$ be the compact subgroup in question. Each
$H_i=H\cap\mathrm{GL}_d(E_i)$ is closed in $H$, because a
finite extension of $\mathbb Q_p$ is complete and therefore
closed in $\overline{\mathbb Q}_p$. Their union is $H$.
The Baire theorem for the compact metric space $H$ implies
that some $H_i$ has nonempty interior. A subgroup with
nonempty interior is open, so its index in $H$ is finite.
Adjoin to $E_i$ the entries of finitely many coset
representatives. The resulting finite extension contains
every matrix in $H$.

For the last assertion, if $R$ is the representation, the map
\[
       g\longmapsto
       \begin{pmatrix}R(g)&z(g)\\0&1\end{pmatrix}               \tag{AV6}
\]
is a continuous representation by (AV1). Its image is compact.
Apply the first assertion to (AV6).
\end{proof}

The local conditions also descend. Suppose that $R$ and $z$
are defined over $E$ and that a local coboundary equation has
a solution after extending coefficients to
$\overline{\mathbb Q}_p$. Choose an $E$-linear map
$\lambda:\overline{\mathbb Q}_p\longrightarrow E$ with
$\lambda(1)=1$, and apply it to the solution's coordinates.
This gives a solution over $E$. At a place above $p$, apply
$\lambda\otimes1$ to the solution in the algebraic tensor
product with $B_{\mathrm{cris}}$. The coefficient action of
the local Galois group is trivial, so this map commutes with
the diagonal action. This proves the same assertion there.
No topological assertion about $\lambda$ is needed: it is
applied to a single vector giving a coboundary, not to a
variable cocycle. Lemma~\ref{av-coefficients} and this observation
justify the use of (AV3) for the cocycles with algebraic
$p$-adic coefficients in the submitted source.

\subsection{Classical Galois representations and vanishing monodromy}
\label{sec-classical-galois}

The Galois representations needed below can be taken directly from
the classical theorems for regular algebraic polarized cuspidal
representations. At places above the coefficient prime, their general
local compatibility theorem is sufficient whenever the automorphic
parameter has zero monodromy. This removes the need for a separate
construction of these global Galois representations.

Let $E$ be a CM field, let $E^+$ be its maximal totally real
subfield, and let $c$ denote complex conjugation. Normalize local
class field theory by sending a uniformizer to geometric Frobenius.
For an irreducible admissible representation $\pi_v$ of
$\mathrm{GL}_n(E_v)$, put
\[
 R_v(\pi)=\operatorname{rec}_{E_v}
       \bigl(\pi_v\otimes|\det|_v^{(1-n)/2}\bigr).
\]
The absolute value is normalized by $|\varpi_v|_v=q_v^{-1}$.
We write $\operatorname{WD}^{\mathrm{F\!\!-ss}}$ for Frobenius
semisimplification of the Weil--Deligne parameter; this operation
retains its monodromy operator.

\begin{theorem}\label{thm-classical-galois-zero-monodromy}
Let $n\geq1$, and let $\pi$ be a regular algebraic cuspidal
automorphic representation of $\mathrm{GL}_n(\mathbb A_E)$.
Suppose that a Hecke character $\chi$ of $E^+$ has $\chi_v(-1)$
independent of the real place $v$ and satisfies
\[
 \pi^c\simeq\pi^\vee\otimes
              (\chi\circ N_{E/E^+}\circ\det).
\]
Fix a prime $\ell$ and a field isomorphism
$\iota:\overline{\mathbb Q}_\ell\simeq\mathbb C$.
Then there are a continuous semisimple representation
\[
 r_{\ell,\iota}(\pi):G_E\longrightarrow
                   \mathrm{GL}_n(\overline{\mathbb Q}_\ell)
\]
and an integer $w$ with the following properties.
For every $v\nmid\ell$,
\[
 \iota\operatorname{WD}
      \bigl(r_{\ell,\iota}(\pi)|_{G_{E_v}}\bigr)^%
          {\mathrm{F\!\!-ss}}
       \simeq R_v(\pi).
\]
At every $v\mid\ell$, the representation is de Rham and
potentially semistable. It is crystalline if $\pi_v$ is spherical.
If the monodromy of $R_v(\pi)$ is zero, the same displayed
local compatibility holds at $v$.

Use the weight convention in which $\pi_\infty$ has the
infinitesimal character of the dual of the algebraic representation
of highest weight
$a_\sigma=(a_{\sigma,1}\geq\cdots\geq a_{\sigma,n})$.
For every embedding $\tau:E\hookrightarrow\overline{\mathbb Q}_\ell$,
\[
 \operatorname{HT}_\tau(r_{\ell,\iota}(\pi))
       =\{a_{\iota\tau,i}+n-i:1\leq i\leq n\},
 \qquad
 \operatorname{HT}_{\tau c}(r_{\ell,\iota}(\pi))
       =\{w-h:h\in\operatorname{HT}_\tau(r_{\ell,\iota}(\pi))\}.
\]
Here the cyclotomic character has Hodge--Tate number $-1$.
\end{theorem}

\begin{proof}
If necessary, multiply $\chi$ by the quadratic character of $E/E^+$.
Its pullback by the norm is trivial, and this arranges
$\chi_v(-1)=(-1)^n$ at every real place without changing $\pi$.
The global representation, compatibility away from $\ell$,
de Rham property, labeled weights, conjugate-weight identity and
crystallinity at spherical places are
\cite[Theorem~2.1.1]{blggt-potential-automorphy}.
Potential semistability follows from
\cite[Theorem~0.6, equivalently Theorem~5.19]{berger-monodromy}.

At $v\mid\ell$, the local theorem
\cite[Theorem~1.1]{blggt-local-global-p} gives the same
semisimplified Weil representation and a monodromy operator
dominated by that of $R_v(\pi)$. Explicitly, within every
inertial class, the partial sums of the largest monodromy block
lengths of the first parameter are at most those of the second.
If the second monodromy is zero, every such block has length one.
The first block length is consequently at most one in every
class, so the first monodromy is zero too. The common
semisimplified Weil representation now gives equality of the
Frobenius-semisimplified parameters. This use of the local theorem
requires neither an Iwahori fixed vector nor an additional
regularity or prime-size condition.
\end{proof}

In particular, the theorem applies at all the toral and tame
places considered here as soon as their specified automorphic
parameters have zero monodromy. The toral and tame restrictions
are unnecessary for this classical Galois assertion. Purity alone
would not give this conclusion: a pure Weil--Deligne
representation can have nonzero monodromy. The argument uses
monodromy dominance and the explicit zero-monodromy hypothesis.

\subsubsection{Weights and changes of realization}

The half-integral shift above can be checked directly. On an
algebraic highest-weight vector of weight
$\nu_0\geq\cdots\geq\nu_{n-1}$, the shifted infinitesimal
parameters are
\[
 z_k=\nu_k+\frac{n-1}{2}-k\qquad(0\leq k<n).
\]
For the dual highest-weight convention,
$\nu_k=-a_{n-k}$. Hence
\[
 \left\{\frac{n-1}{2}-z_k:0\leq k<n\right\}
       =\{a_i+n-i:1\leq i\leq n\}.
\]
The de Rham filtration gives the same sign: for the cyclotomic
character its invariant vector has the form $v\otimes t^{-1}$,
whose filtration degree is $-1$.

All comparisons in the next lemma use the induced identifications of
coefficient fields and embedding labels.

\begin{lemma}\label{lem-galois-period-realizations}
The crystalline, de Rham and potentially semistable conditions,
the labeled Hodge--Tate multiplicities and the Weil--Deligne
isomorphism class are unchanged by enlarging the finite
coefficient field, changing compatible algebraic closures and
completion embeddings, or changing presentations of the period
rings by equivariant isomorphisms preserving their scalar
embeddings, Frobenius, filtration and logarithmic period.
These changes may be made separately at every finite place
after one global representation has been fixed.
\end{lemma}

\begin{proof}
Write $L'/L$ for a finite extension of coefficient fields.
Choosing an $L$-basis of $L'$ shows, coordinate by coordinate,
that tensoring with $L'$ commutes with taking invariants.
The same argument applies to every filtered embedding
eigenspace, so it preserves the dimensions of its successive
quotients. Potentially semistable periods are the union of
semistable invariant spaces over finite extensions of the local
field. Finitely many vectors in this union are already defined
over one common finite extension; consequently coefficient
extension commutes with this construction as well. Frobenius
and monodromy commute with these identifications.

An equivariant period-ring isomorphism, tensored with the
identity of the representation, identifies the same invariant
spaces and filtered eigenspaces. In the usual presentation
$B_{\mathrm{st}}=B_{\mathrm{cris}}[X]$ the operators are
\[
 gX=X+c(g)t,\qquad \varphi X=\ell X,
 \qquad N=-\frac{d}{dX}.
\]
Changing the compatible roots of $\ell$ replaces $X$ by
$X+bt$ for some $b\in\mathbb Z_\ell$. Changing the primitive
roots of unity scales $t$ and inversely scales $c(g)$.
Both changes give the same equivariant period construction.

If the residue field has size $q=\ell^f$ and $g$ has
arithmetic degree $d$, the associated Weil action is
$g\varphi^{-fd}$. For geometric degree $m=-d$ this is
$g\varphi^{fm}$. The identity $N\varphi=\ell\varphi N$
therefore gives
$r(g)Nr(g)^{-1}=q^{-m}N$, as required by the normalization
fixed above. Scalar specialization and Frobenius
semisimplification preserve all these isomorphisms.

Finally, Lemma~\ref{av-coefficients} provides a finite
coefficient model of the given global representation. It can
be enlarged to contain any further finite local coefficient
field. Thus all the local choices can follow the single
choice of the global representation and its labeled weights.
\end{proof}

The later adjoint theorem is applied under its stated irreducibility
hypotheses.

\subsection{The polarization and its adjoint action}

Let now $F$ be a CM field with maximal totally real subfield
$F^+$. Fix a complex conjugation $c\in G_{F^+}$, so that
$c^2=1$ and $c\notin G_F$. Let $n\ge2$, let
$\varepsilon_p:G_{F^+}\longrightarrow\mathbb Z_p^\times$
be the cyclotomic character, and fix an abstract field
isomorphism $\iota:\overline{\mathbb Q}_p\simeq\mathbb C$.
This isomorphism is not asserted to be continuous.

For a representation
$r:G_F\longrightarrow\mathrm{GL}_n(\overline{\mathbb Q}_p)$,
write $r^c(g)=r(cgc^{-1})$. A matrix $P\in
\mathrm{GL}_n(\overline{\mathbb Q}_p)$ gives the required
polarization when
\[
           r(g)^t P r^c(g)=\varepsilon_p(g)^{1-n}P
                         \qquad(g\in G_F).                    \tag{AV7}
\]
For a symmetric such $P$, the polarized adjoint representation
$\operatorname{Ad}_P(r)$ of $G_{F^+}$ acts on
$M_n(\overline{\mathbb Q}_p)$ by
\[
 \begin{aligned}
 g\cdot X&=r(g)Xr(g)^{-1}&& (g\in G_F),\\
 c\cdot X&=-P^{-1}X^tP.
 \end{aligned}                                                 \tag{AV8}
\]
The square of the second operator is the identity. Transposing
(AV7), or solving it for $r^c(g)$, verifies that this operator
conjugates $\operatorname{Ad}(r(g))$ to
$\operatorname{Ad}(r^c(g))$. These are the defining relations
for the extension from the subgroup $G_F$ of index two.

\begin{theorem}[Classical adjoint vanishing]\label{av-thorne}
Suppose that $\pi$ is a regular algebraic cuspidal automorphic
representation of $\mathrm{GL}_n(\mathbb A_F)$ satisfying
$\pi^c\simeq\pi^\vee$. Let $r_{\pi,\iota}$ be its continuous
semisimple Galois representation, in the Tate normalization.
If
\[
 r_{\pi,\iota}|_{G_{F(\mu_{p^\infty})}}
                         \quad\hbox{is irreducible},            \tag{AV9}
\]
then it has a symmetric polarization as in (AV7), and
\[
 H^1_f\bigl(F^+,\operatorname{Ad}_P(r_{\pi,\iota})\bigr)=0.
                                                               \tag{AV10}
\]
This holds for every prime $p$; no condition on its residual
representation is imposed.
\end{theorem}

The vanishing is Thorne's Theorem~6.1, specialized to the
conjugate-self-dual case. The symmetric pairing and the
minus sign in (AV8) are the conventions in Newton--Thorne,
page~1920; the Frobenius normalization is fixed on page~1923.
\cite{thorne-adjoint-2022,newton-thorne-adjoint-2023}
In the submitted source's convention a polarized pair also
records a character $\chi$ of $F^+$ satisfying
$\chi_v(-1)=(-1)^n$. For the above $\pi$ take
$\chi=\delta_{F/F^+}^n$, where $\delta_{F/F^+}$ is the
quadratic Hecke character of the CM extension. Its pullback
by the idele norm from $F$ is trivial, and its component at
each real place is the sign character. Its Galois character
restricts trivially to $G_F$. This explains both the parity
condition and the multiplier $\varepsilon_p^{1-n}$ in (AV7).
In odd rank one must not replace this auxiliary character by
the trivial character in that particular convention.

\begin{lemma}\label{av-pairing-unique}
If $r$ is irreducible and $P_0,P$ both satisfy (AV7), then
$P$ is a nonzero scalar multiple of $P_0$. Consequently the
operator in the second line of (AV8) is independent of the
choice of polarization.
\end{lemma}
\begin{proof}
Substitution of $P=P_0T$ in (AV7) shows that $T$ commutes
with every $r^c(g)$. An eigenspace of $T$ exists over
$\overline{\mathbb Q}_p$, and is stable under $r^c$.
Irreducibility forces this eigenspace to be the whole space,
so $T$ is scalar. Invertibility makes the scalar nonzero,
and it cancels in $P^{-1}X^tP$.
\end{proof}

Scalar twists cause no difficulty either. For a character
$\alpha$ of $G_F$, conjugation by $\alpha(g)r(g)$ equals
conjugation by $r(g)$. If the polarized extension uses the
same pairing, its second operator in (AV8) is identical too.
In particular scalar twisting changes neither the invariant
subspaces in (AV9) nor this adjoint representation.

\subsection{Comparison with the actual automorphic point}

We record explicitly the elementary uniqueness principle used
when the Galois representation is specified by Frobenius
polynomials. If two continuous semisimple representations
over $\overline{\mathbb Q}_p$ have equal characteristic
polynomials at all but finitely many unramified Frobenius
classes, they are isomorphic. Lemma~\ref{av-coefficients}
puts both over a finite extension, and Chebotarev density
and continuity give equality of their traces at every group
element. To see the last algebraic implication, let $A$ be
the finite-dimensional algebra spanned by the image on the
direct sum of the two representation spaces. The direct
sum is a faithful semisimple $A$-module, so the radical of
$A$ is zero. Over an algebraically closed field, $A$ is a
product of matrix algebras. The trace of each central
primitive idempotent gives the multiplicity of the
corresponding simple module, multiplied by its nonzero
dimension. Equality of traces on group elements implies
equality on their span, hence on these idempotents. The
multiplicities agree, proving the assertion. This is the
usual Brauer--Nesbitt argument, included to specify exactly
what the Frobenius comparison supplies.

Here is the quotient form needed in the source. Let
$q:G_{F^+}\twoheadrightarrow\Gamma$ be a continuous quotient
of profinite groups. Suppose that a quadratic character
$\eta:\Gamma\longrightarrow\{\pm1\}$ and a continuous
character $\epsilon:\Gamma\longrightarrow\mathbb Z_p^\times$
satisfy
\[
 \ker(\eta q)=G_F,\qquad
 \epsilon q=\varepsilon_p.
                                                               \tag{AV11}
\]
Put $\Gamma_0=\ker\eta$, $c_0=q(c)$ and
$\Gamma_1=\ker(\epsilon|_{\Gamma_0})$.

\begin{proposition}\label{av-quotient}
Let $\rho:\Gamma_0\longrightarrow
\mathrm{GL}_n(\overline{\mathbb Q}_p)$ be continuous, and
assume that $\rho|_{\Gamma_1}$ is irreducible. Suppose that
$r=\rho\circ q|_{G_F}$ is attached to a regular algebraic
conjugate-self-dual cuspidal representation as in
Theorem~\ref{av-thorne}, by the Frobenius criterion just
described. Let $P$ be any matrix satisfying the quotient
version of (AV7), and let $R$ be a representation of
$\Gamma$ satisfying
\[
 R(g)X=\rho(g)X\rho(g)^{-1}\quad(g\in\Gamma_0),\qquad
 R(c_0)X=-P^{-1}X^tP.                                          \tag{AV12}
\]
If $z:\Gamma\longrightarrow M_n(\overline{\mathbb Q}_p)$
is a continuous one-cocycle whose pullback along $q$
satisfies (AV2) at every finite place of $F^+$, then there is
one matrix $X$ such that
\[
                          z(g)=R(g)X-X\qquad(g\in\Gamma).
                                                               \tag{AV13}
\]
\end{proposition}
\begin{proof}
Surjectivity of $q$ and the first equality in (AV11) give
$q(G_F)=\Gamma_0$. The second equality then gives the exact
identity
\[
          q\bigl(G_{F(\mu_{p^\infty})}\bigr)=\Gamma_1.
                                                               \tag{AV14}
\]
Indeed, lift an element of $\Gamma_1$ to $G_{F^+}$;
the first equality forces the lift into $G_F$, and the
second forces its cyclotomic character to be one. The
reverse inclusion follows immediately from (AV11).
Thus the assumed irreducibility is precisely (AV9).

The Frobenius criterion identifies $r$ with $r_{\pi,\iota}$.
Pulling $P$ back gives (AV7). By
Lemma~\ref{av-pairing-unique}, it differs by a scalar from
the symmetric pairing of Theorem~\ref{av-thorne}. Hence
$R\circ q$ is its polarized adjoint representation: the
two actions agree on $G_F$ and on $c$, which together
generate $G_{F^+}$. Notice that the proposition did not
need a separate symmetry assumption on $P$.

The representation $R$ is continuous. This follows from
(AV12), continuity of $\rho$, and the two open cosets of
$\Gamma_0$. By Lemma~\ref{av-coefficients} the cocycle
and representation have a common finite coefficient field;
the local descent argument following that lemma applies.
Theorem~\ref{av-thorne} therefore gives
\[
                         z(q(\sigma))=R(q(\sigma))X-X
                         \qquad(\sigma\in G_{F^+}).
\]
Surjectivity of $q$ proves (AV13).
\end{proof}

\begin{corollary}\label{av-source}
The adjoint-vanishing assertion at the automorphic points of
the submitted ordinary deformation construction follows from
the classical theorem above.
\cite{alpo-uniform-source}
\end{corollary}
\begin{proof}
Use $F=L$, $F^+=L^+$, $\Gamma=G_{L^+,S'}$ and the
natural quotient $q$. The source's level data state (AV11)
exactly, and its selected $c_0$ is the image of a genuine
complex conjugation. The irreducibility assumption concerns
exactly the subgroup $\Gamma_1$. The automorphic-point
hypothesis supplies the cuspidal representation, its regular algebraicity,
conjugate self-duality, and the matching semisimple Galois
representation. In the source's Frobenius convention the
polynomial of $r(\sigma^{-1})$, for arithmetic Frobenius
$\sigma$, has Satake roots multiplied by
$q_v^{(n-1)/2}$. This is the same Tate normalization used
in Theorem~\ref{av-thorne}.

The specified polarization and adjoint action are precisely
(AV7) and (AV12). The local condition says that the
restriction of the cocycle is a coboundary on inertia away
from $p$, and becomes one after tensoring with
$B_{\mathrm{cris}}$ above $p$. It initially chooses one
embedding above each place; (AV4) shows independence of
that choice. Its local group consists of pairs of global
and local automorphisms which agree on the chosen embedded
algebraic closure of the number field. Projection to the
local factor identifies it with the Galois group of the
corresponding completion: the local automorphism has a
unique restriction to that algebraic closure, and fixing
the number field is equivalent, by continuity, to fixing
its completion. Thus these are exactly (AV2).
Proposition~\ref{av-quotient}
therefore gives the requested global coboundary conclusion.
All extra ordinary, weight, rank-at-most-four and local
genericity assumptions in the source implication may be
retained, but are unnecessary for this deduction.
\end{proof}

\subsection{The symplectic projection}

For later use we also give the matrix identity which transports
this vanishing to symplectic infinitesimal deformations.
Let $J\in\mathrm{GL}_{2m}(E)$ satisfy $J^t=-J$, and put
\[
 X^*=J^{-1}X^tJ,\qquad
 \mathfrak{sp}(J)=\{X:X^*=-X\},\qquad
                        P_-(X)=\frac{X-X^*}{2}.               \tag{AV15}
\]
The symbol $P_-$ denotes a linear projection and is unrelated
to the polarization matrix $P$ above. The coefficient field
has characteristic zero, including when $p=2$.
Suppose that $r:G\longrightarrow\mathrm{GSp}(J,E)$ is a
representation and $\eta:G\longrightarrow\{\pm1\}$ is
a character. If $A_g(X)=r(g)Xr(g)^{-1}$, then
\[
 R_g(X)=A_g(P_-X)+\eta(g)A_g(X-P_-X)                            \tag{AV16}
\]
defines a representation. To check this, the similitude
identity gives $(A_gX)^*=A_g(X^*)$, so $P_-$ commutes
with every $A_g$. It is idempotent, and (AV16) acts by
$A_g$ on its image and by $\eta(g)A_g$ on its kernel.
The group law follows on these two summands.

When $\eta(c)=-1$, (AV16) gives
\[
 R_c(X)=-r(c)J^{-1}X^tJr(c)^{-1}
       =-(Jr(c)^{-1})^{-1}X^t(Jr(c)^{-1}).                      \tag{AV17}
\]
For $g\in\ker\eta$, its action is the ordinary adjoint
action. Thus (AV17) is exactly the polarized extension
with pairing matrix $Jr(c)^{-1}$, whenever the latter
representation has the automorphic polarization considered
above.

The inclusion $\mathfrak{sp}(J)\hookrightarrow M_{2m}(E)$
and projection $P_-$ are equivariant for these actions.
They commute with restriction to inertia and with the
diagonal action after tensoring with $B_{\mathrm{cris}}$.
They therefore give split maps on the Selmer groups in
(AV3). In particular vanishing for $R$ implies vanishing
for its symplectic summand. Explicitly, if a cocycle
$z(g)\in\mathfrak{sp}(J)$ becomes $R_gX-X$, then
\[
                    z(g)=A_g(P_-X)-P_-X.                       \tag{AV18}
\]
This retains the primitive in the symplectic Lie algebra;
it is stronger than merely finding a primitive in all
matrices. The argument is the ordinary matrix calculation
needed in the source's rank-four rigidity step.

\paragraph{What remains.}
Corollary~\ref{av-source} closes the adjoint-vanishing
implication at the actual automorphic points. To deduce
the remaining complementary-parity estimates, one still
needs the arithmetic construction and automorphy of the
deformations which encode the original rank-one Selmer
classes, and the quantitative integral comparison with
the relevant congruence modules. Rational vanishing does
not bound the length of a torsion module: the modules
$\mathbb Z_p/p^a\mathbb Z_p$ all become zero after
tensoring with $\mathbb Q_p$, while their lengths are
the arbitrary integers $a$. Thus neither the polynomial
exponent nor the coefficient of the common interval
bound $B$ can be inferred from (AV10) alone. These
remaining steps must be checked with their original
uniform quantifiers.


